%=============================================================================!
% 2.8  THE SPRING                                                             !
%=============================================================================!
\section{The spring}
\label{sec:ne-spring}

A block fixed to a spring and pulled aside swings back and forth about
its resting place. This motion recurs throughout the book because an atom
displaced slightly from a stable position feels a similar restoring
force from its neighbours. Chapter~4 develops that connection; here we derive the
motion from the second law.

\subsection*{Hooke's law}

Lay a spring on a smooth table, fix one end to a wall and attach a block
of mass $m$ to the other. We measure the position $x$ of the block from
its resting place, where the spring has its natural length, so that $x$
is the stretch of the spring: positive when it is stretched, negative
when it is squashed. A stretched spring pulls the block back towards the
wall, and a squashed one pushes it away (\cref{fig:ne-spring-states}).
Measurements with a spring balance show that, for stretches small
compared with the length of the spring, the force is proportional to the
stretch,
\begin{equation*}
   F = -kx .
\end{equation*}
This is Hooke's law. The stiffness $k$, in newtons per metre, measures how
hard the spring is to stretch, and the minus sign records that the force
always points back towards $x = 0$, against the stretch. A spring with $k
= \SI{50}{\newton\per\metre}$ needs a force of \SI{1}{\newton} to stretch
it by $1/50 = \SI{0.02}{\metre}$, or \SI{2}{\centi\metre}.

\begin{figure}[htb]
\centering
\begin{tikzpicture}[line cap=round,
   force/.style={-{Stealth[length=5pt]}, line width=1.1pt, accent},
   stretch/.style={{Stealth[length=3pt]}-{Stealth[length=3pt]},
      line width=0.4pt},
   coil/.style={decorate, decoration={coil, aspect=0.45,
      segment length=4pt, amplitude=3.5pt, pre length=3pt,
      post length=3pt}, line width=0.6pt}]
   % one row per state; \x is the stretch, in the units of the drawing
   \foreach \row/\x/\name in {0/-0.6/squashed, 1/0/natural length,
      2/0.8/stretched} {
      \begin{scope}[yshift=-\row*1.9cm]
         \fill[black!35] (-0.12,-0.35) rectangle (0,0.45);
         \draw[line width=0.6pt] (0,-0.35) -- (6.0,-0.35);
         \draw[coil] (0,0.05) -- (3.0+\x,0.05);
         \draw[line width=0.6pt, fill=black!8] (3.0+\x,-0.35)
            rectangle (3.8+\x,0.45);
         \node[anchor=west] at (6.3,0.05) {\small \name};
      \end{scope}
   }
   % the resting place of the block's left face
   \draw[dashed, black!50, line width=0.5pt] (3.0,1.3) -- (3.0,-4.6);
   \node[above] at (3.0,1.3) {\small $x = 0$};
   % the force of the spring on the block, drawn above it
   \draw[force] (2.6,0.8) -- (3.5,0.8) node[right] {\small $F$};
   \draw[force] (4.7,-3.0) -- (3.8,-3.0) node[left] {\small $F$};
   % the stretch, measured from the resting place
   \draw[stretch] (3.0,-0.65) -- (2.4,-0.65)
      node[left] {\scriptsize $x < 0$};
   \draw[stretch] (3.0,-4.45) -- (3.8,-4.45)
      node[right] {\scriptsize $x > 0$};
\end{tikzpicture}
\caption{A block on a spring. Squashed ($x < 0$), the spring pushes the
block away from the wall; at its natural length ($x = 0$) it exerts no
force; stretched ($x > 0$), it pulls the block back. In each case the
force (teal) points towards $x = 0$, against the stretch.}
\label{fig:ne-spring-states}
\end{figure}

\subsection*{The equation of motion}

Vertically, the weight of the block is balanced by the push of the table,
called the normal force because it acts at right angles to the surface,
which mathematicians call normal to it. Along the table the spring is the
only force, so the second law reads $m\,\dd^2 x/\dd t^2 = -kx$, and
dividing by $m$,
\begin{equation}
   \frac{\dd^2 x}{\dd t^2} = -\frac{k}{m}\,x .
   \label{eq:ne-spring}
\end{equation}
Rather than solve this equation from nothing, we recognise it. In
\cref{sec:mo-trig} we found that every sinusoid obeys $\dd^2 x/\dd t^2 =
-\omega^2 x$, which is \cref{eq:ne-spring} if $\omega^2 = k/m$. With
$\omega = \sqrt{k/m}$, both $\cos\omega t$ and $\sin\omega t$ are such
sinusoids, and the equation is obeyed also by any combination
\begin{equation*}
   x(t) = c_1\cos\omega t + c_2\sin\omega t
\end{equation*}
with constant $c_1$ and $c_2$. Differentiating twice with the chain rule
confirms it,
\begin{align*}
   \frac{\dd x}{\dd t} &= -c_1\omega\sin\omega t + c_2\omega\cos\omega t ,\\
   \frac{\dd^2 x}{\dd t^2} &= -c_1\omega^2\cos\omega t
      - c_2\omega^2\sin\omega t
      && \text{differentiating again}\\
   &= -\omega^2 x = -\frac{k}{m}\,x
      && \text{taking out the common factor $-\omega^2$,}
\end{align*}
so \cref{eq:ne-spring} holds whatever the values of $c_1$ and $c_2$.

\subsection*{Matching the starting state}

The two constants are fixed by the state at $t = 0$. Since $\cos 0 = 1$
and $\sin 0 = 0$, setting $t = 0$ in $x$ and $\dd x/\dd t$ gives $x_0 =
c_1$ and $v_0 = c_2\omega$, where $x_0$ and $v_0$ are the starting position
and velocity. The motion is therefore
\begin{equation}
   x(t) = x_0\cos\omega t + \frac{v_0}{\omega}\,\sin\omega t .
   \label{eq:ne-spring-solution}
\end{equation}
Every starting state $(x_0, v_0)$ is matched in this way, so by the
uniqueness of solutions (\cref{sec:ne-equations})
\cref{eq:ne-spring-solution} is the motion for every start. With $x_0 =
v_0 = 0$ it gives $x = 0$ for all time: a block at rest at its resting
place stays there.

The same motion can be written as a single sinusoid, which shows at a
glance how far the block swings. We write it with a cosine,
\begin{equation*}
   x(t) = A\cos(\omega t + \phi) ,
\end{equation*}
which by the quarter-turn relation $\cos\theta = \sin(\theta + \pi/2)$ is
the sinusoid of \cref{eq:mo-sinusoid} with its phase larger by $\pi/2$.
Expanding it with the addition formula \cref{eq:mo-addition},
\begin{equation*}
   A\cos(\omega t + \phi) = A\cos\phi\,\cos\omega t - A\sin\phi\,\sin\omega t ,
\end{equation*}
which is \cref{eq:ne-spring-solution} when the coefficients of $\cos\omega
t$ and $\sin\omega t$ agree,
\begin{equation*}
   A\cos\phi = x_0, \qquad A\sin\phi = -\frac{v_0}{\omega} .
\end{equation*}
To find $A$, we square both equations and add them, using $\cos^2\phi +
\sin^2\phi = 1$:
\begin{equation*}
   x_0^2 + \Bigl(\frac{v_0}{\omega}\Bigr)^2
   = A^2\bigl(\cos^2\phi + \sin^2\phi\bigr) = A^2
   \quad\Longrightarrow\quad
   A = \sqrt{x_0^2 + (v_0/\omega)^2} .
\end{equation*}
If $A = 0$ the block rests at $x = 0$ and the phase plays no part.
Otherwise the two equations give $\cos\phi = x_0/A$ and $\sin\phi =
-v_0/(\omega A)$. Together these fix the point $(\cos\phi, \sin\phi)$ on
the unit circle, and so fix $\phi$ up to whole turns of $2\pi$. The
arcsine of \cref{sec:mo-trig} returns the angle between $-\pi/2$ and
$\pi/2$ with the right sine; the cosine of that angle is never negative,
and by $\cos^2 + \sin^2 = 1$ it equals $+\sqrt{1 - \sin^2\phi}$, which is
$\cos\phi$ itself when $\cos\phi \ge 0$. In that case, therefore, $\phi =
\arcsin(\sin\phi)$.
When $\cos\phi < 0$ the angle is $\pi - \arcsin(\sin\phi)$ instead, which
has the same sine, since $\sin(\pi - \theta) = \sin\theta$, and a
negative cosine, since $\cos(\pi - \theta) = -\cos\theta$. The sign of
the cosine thus selects between the two angles that share a sine.

Take a block of \SI{0.5}{\kilogram} on the spring of
\SI{50}{\newton\per\metre}, so that $\omega = \sqrt{50/0.5} =
\SI{10}{\radian\per\second}$, and start it \SI{3}{\centi\metre} from its
resting place, moving outwards, away from the wall, at
\SI{0.4}{\metre\per\second}. In metres,
$x_0 = 0.03$ and $v_0/\omega = 0.04$, so $x(t) = 0.03\cos 10t +
0.04\sin 10t$ by \cref{eq:ne-spring-solution}, and
\begin{align*}
   A &= \sqrt{0.03^2 + 0.04^2} = \SI{0.05}{\metre} ,\\
   \cos\phi &= 0.03/0.05 = 0.6, \qquad
   \sin\phi = -0.04/0.05 = -0.8 .
\end{align*}
The block swings out to \SI{5}{\centi\metre} before turning back. The
cosine is positive, so $\phi = \arcsin(-0.8) = \SI{-0.927}{\radian}$.
Started instead \SI{3}{\centi\metre} on the wall side, $x_0 = -0.03$,
with the same velocity, the block has the same amplitude, but $\cos\phi =
-0.6$ is negative, so $\phi = \pi - \arcsin(-0.8) = \pi + 0.927 =
\SI{4.069}{\radian}$.

\subsection*{The period and the motion}

The block repeats its motion when $\omega t$ grows by $2\pi$, so its
angular frequency and period are
\begin{equation}
   \omega = \sqrt{\frac{k}{m}}, \qquad
   \frac{2\pi}{\omega} = 2\pi\sqrt{\frac{m}{k}} .
   \label{eq:ne-spring-omega}
\end{equation}
A heavier block oscillates more slowly and a stiffer spring more quickly.
As long as Hooke's law holds, the period does not depend on the
amplitude, and there is a simple reason.
Pull the block twice as far and the spring pulls twice as hard, so at
every stage the acceleration is twice as large; the block covers twice the
distance at twice the speed and takes the same time. In the language of
the equation, if $x(t)$ solves \cref{eq:ne-spring} then so does $2x(t)$,
because both sides double.

\Cref{fig:ne-spring} shows the simplest start of all, the same block
pulled out \SI{4}{\centi\metre} and released from rest, for which $\phi
= 0$, since then $x_0 = A$ and $v_0 = 0$. The period is $2\pi/10 =
\SI{0.628}{\second}$. On the way in, the spring pulls the block in the
direction it is moving, so the block speeds up all the way to the centre;
past the centre the spring pulls against the motion and slows it. The
speed is therefore largest at $x = 0$. Differentiating $x = A\cos(\omega
t + \phi)$ with the chain rule gives $\dd x/\dd t = -A\omega\sin(\omega t
+ \phi)$; where $x = 0$ the cosine is zero, so the sine is $\pm1$ by
$\cos^2 + \sin^2 = 1$, and the speed there is $A\omega = 0.04\times 10 =
\SI{0.4}{\metre\per\second}$. At the turning points, $x = \pm A$, the
block is momentarily at rest, yet its acceleration there, $-\omega^2 x =
\mp A\omega^2 = \mp\SI{4}{\metre\per\second\squared}$ by
\cref{eq:ne-spring}, is the largest it ever has, since $\lvert x\rvert$
never exceeds $A$. This is the top of the throw of \cref{sec:mo-acceleration} again:
zero velocity does not mean zero acceleration.

\begin{figure}[htb]
\centering
\includegraphics{ch02_newton/spring}
\caption{A block of \SI{0.5}{\kilogram} on a spring of
\SI{50}{\newton\per\metre}, released from rest \SI{4}{\centi\metre} from
its resting place, over two periods. (a) Position, (b) velocity and (c)
acceleration. At the dotted lines the block turns: its velocity is zero
and its acceleration largest. At the dashed lines it passes the centre:
its speed is largest and its acceleration zero.}
\label{fig:ne-spring}
\end{figure}

\notebook{02}{change the mass, stiffness and starting state of the block}

\begin{selfcheck}
A block oscillates on a spring with a period of \SI{1}{\second}. What is
the period if the mass is made four times larger? What if the block is
pulled out twice as far before release?
\end{selfcheck}
